\chapter{Triangles and Angles}\label{ch:g7:triangles}

Can you draw a triangle with sides $10$, $3$ and $2$? With angles
$80^\circ$, $70^\circ$ and $50^\circ$? This chapter answers both
questions with two of the most useful facts of plane geometry: the
triangle inequality and the $180^\circ$ angle \omterm{def:g1:addition:def}{sum} --- plus the angle
vocabulary that comes with crossing and \omterm{def:g6:lines:perp}{parallel lines}.

\section{Angles around crossing lines}

\begin{definition}[Angle pairs]\label{def:g7:triangles:pairs}
\begin{itemize}
  \item Two angles are \emph{complementary} when their measures add up
        to $90^\circ$, \emph{supplementary} when they add up to
        $180^\circ$.
  \item When two lines cross, the angles facing each other across the
        crossing point are \emph{vertically
        opposite}\index{vertically opposite angles}.
  \item When a line crosses two other lines, angles located in the
        same position at each crossing are
        \emph{corresponding angles}; angles between the two lines on
        opposite sides of the crossing line are \emph{alternate
        angles}\index{alternate angles}.
\end{itemize}
\end{definition}

\begin{theorem}[Angle equalities]\label{thm:g7:triangles:equalities}
\begin{enumerate}
  \item \omterm{def:g7:triangles:pairs}{Vertically opposite angles} are equal.
  \item If two lines are \omterm{def:g6:lines:perp}{parallel}, corresponding angles are equal, and
        \omterm{def:g7:triangles:pairs}{alternate angles} are equal. Conversely, equal corresponding
        (or alternate) angles force the lines to be \omterm{def:g6:lines:perp}{parallel}.
\end{enumerate}
\end{theorem}

\begin{proof}[Proof of 1]
The two angles $\widehat{1}$ and $\widehat{2}$ on one side of a line
make a straight angle: $\widehat 1 + \widehat 2 = 180^\circ$. The
\omterm{def:g7:triangles:pairs}{vertically opposite} angle $\widehat 3$ of $\widehat 1$ also satisfies
$\widehat 2 + \widehat 3 = 180^\circ$ (same reason). So
$\widehat 1 = 180^\circ - \widehat 2 = \widehat 3$. Point 2 is
admitted at this level.
\end{proof}

\begin{omfigure}
\begin{tikzpicture}[scale=1.0]
  \draw[thick] (-2.0,-0.9) -- (2.0,0.9);
  \draw[thick] (-2.0,0.9) -- (2.0,-0.9);
  \draw[omDef, ->] (0.75,0) arc (0:24.2:0.75);
  \draw[omDef, ->] (-0.75,0) arc (180:204.2:0.75);
  \node[omDef, font=\small] at (1.15,0.28) {$\widehat 1$};
  \node[omDef, font=\small] at (-1.15,-0.28) {$\widehat 3$};
  \draw[omThm, ->] (0.55,0.25) arc (24.2:155.8:0.6);
  \node[omThm, font=\small] at (0,0.75) {$\widehat 2$};
  \begin{scope}[xshift=7.2cm]
  \draw[thick] (-2.2,-1.0) -- (2.2,-1.0);
  \draw[thick] (-2.2,0.8) -- (2.2,0.8);
  \draw[thick, omExo] (-1.6,-1.6) -- (1.6,1.5);
  \draw[omDef, ->] (-0.28,-1.0) ++(0.75,0) arc (0:44:0.75);
  \node[omDef, font=\small] at (0.85,-0.62) {$\widehat a$};
  \draw[omDef, ->] (0.63,0.8) ++(0.75,0) arc (0:44:0.75);
  \node[omDef, font=\small] at (1.76,1.18) {$\widehat b$};
  \end{scope}
\end{tikzpicture}

{\small Left: \omterm{def:g7:triangles:pairs}{vertically opposite angles} $\widehat 1 = \widehat 3$
(each is supplementary to $\widehat 2$). Right: \omterm{def:g6:lines:perp}{parallel lines} cut by
a third line --- the corresponding angles $\widehat a$ and $\widehat b$
are equal.}
\end{omfigure}

\section{The sum of the angles of a triangle}

\begin{theorem}[Angle sum]\label{thm:g7:triangles:sum}
In every triangle, the three angles add up to $180^\circ$.
\end{theorem}

\begin{proof}
Let $ABC$ be a triangle. Draw the line through $A$ \omterm{def:g6:lines:perp}{parallel} to
$(BC)$. At the \omterm{def:g5:solids:def}{vertex} $A$, three angles line up along this \omterm{def:g6:lines:perp}{parallel}
and make a straight angle, $180^\circ$: the middle one is
$\widehat{A}$, and the two outer ones are \omterm{def:g7:triangles:pairs}{alternate angles} with
$\widehat B$ and $\widehat C$ (\omterm{def:g6:lines:perp}{parallel lines}!), so they equal
$\widehat B$ and $\widehat C$. Hence
$\widehat B + \widehat A + \widehat C = 180^\circ$.
\end{proof}

\begin{omfigure}
\begin{tikzpicture}[scale=1.15]
  \coordinate (B) at (0,0);
  \coordinate (C) at (4.2,0);
  \coordinate (A) at (1.6,2.0);
  \draw[thick] (B) -- (C) -- (A) -- cycle;
  \draw[thick, omExo] (-0.3,2.0) -- (3.9,2.0);
  \node[above, font=\small] at (A) {$A$};
  \node[below left, font=\small] at (B) {$B$};
  \node[below right, font=\small] at (C) {$C$};
  \draw[omDef, ->] (0.55,0) arc (0:51.3:0.55);
  \node[omDef, font=\small] at (0.95,0.32) {$\widehat B$};
  \draw[omMeth, ->] (3.68,0) arc (180:142.4:0.52);
  \node[omMeth, font=\small] at (3.25,0.3) {$\widehat C$};
  \draw[omDef, ->] (A) ++(-0.55,0) arc (180:231.3:0.55);
  \node[omDef, font=\small] at (0.85,1.75) {$\widehat B$};
  \draw[omMeth, ->] (A) ++(0.52,0) arc (0:-37.6:0.52);
  \node[omMeth, font=\small] at (2.35,1.78) {$\widehat C$};
\end{tikzpicture}

{\small The proof in one picture: along the \omterm{def:g6:lines:perp}{parallel} to $(BC)$ through
$A$, the angles $\widehat B$, $\widehat A$, $\widehat C$ (\omterm{def:g7:triangles:pairs}{alternate
angles} in matching colors) fill a straight angle.}
\end{omfigure}

\begin{example}\label{ex:g7:triangles:sum}
Two angles of a triangle measure $67^\circ$ and $58^\circ$. The third
measures
\[
180 - (67 + 58) = 180 - 125 = 55^\circ .
\]
Special cases worth memorizing: each angle of an \omterm{def:g6:shapes:triangles}{equilateral triangle}
is $180 \div 3 = 60^\circ$; in a \omterm{def:g6:shapes:triangles}{right triangle} the two acute angles
are complementary (they add up to $180 - 90 = 90^\circ$); the base
angles of an \omterm{def:g6:shapes:triangles}{isosceles triangle} are equal, so each is
$\frac{180 - \text{apex}}{2}$.
\end{example}

\begin{example}[No triangle with two right angles]\label{ex:g7:triangles:tworight}
A triangle with two angles of $90^\circ$ would already use up
$180^\circ$, leaving $0^\circ$ for the third --- impossible. The angle
\omterm{def:g1:addition:def}{sum} forbids many triangles at once.
\end{example}

\section{The triangle inequality}

\begin{theorem}[Triangle inequality]\label{thm:g7:triangles:inequality}
In every triangle, each side is shorter than the \omterm{def:g1:addition:def}{sum} of the two
others: for a triangle with sides $a$, $b$, $c$,
\[
a < b + c .
\]
Conversely, three lengths with the \emph{largest} smaller than the \omterm{def:g1:addition:def}{sum}
of the two others can always be made into a triangle. If the largest
\emph{equals} the \omterm{def:g1:addition:def}{sum} of the others, the ``triangle'' is flat: the
three \omterm{def:g5:solids:def}{vertices} are aligned.
\end{theorem}
\admitted

\begin{example}\label{ex:g7:triangles:inequality}
Can sides be $10$, $3$ and $2$? Test the largest: is $10 < 3 + 2 = 5$?
No --- no such triangle exists. The compass shows why: arcs of radii
$3$ and $2$ drawn from the two ends of a \omterm{def:g6:lines:objects}{segment} of length $10$ never
meet.

Sides $7$, $5$ and $4$? Largest: $7 < 5 + 4 = 9$: yes, the triangle
exists (checking the largest side is enough).
\end{example}

\begin{omfigure}
\begin{tikzpicture}[scale=0.72]
  \draw[thick] (0,0) -- (10*0.72,0);
  \fill (0,0) circle (1.5pt) node[below, font=\small] {$A$};
  \fill (7.2,0) circle (1.5pt) node[below, font=\small] {$B$};
  \draw[omDef, thick] (2.16,0) arc (0:60:2.16);
  \draw[omDef, thick] (2.16,0) arc (0:-30:2.16);
  \draw[omThm, thick] (7.2,0) ++(-1.44,0) arc (180:120:1.44);
  \draw[omThm, thick] (7.2,0) ++(-1.44,0) arc (180:210:1.44);
  \node[omDef, font=\small] at (1.1,1.6) {radius $3$};
  \node[omThm, font=\small] at (6.4,1.3) {radius $2$};
\end{tikzpicture}

{\small Trying to build a triangle with sides $10$, $3$, $2$: the two
compass arcs stay hopelessly far apart, because $3 + 2 < 10$.}
\end{omfigure}

\begin{method}[Existence and construction of a triangle]\label{met:g7:triangles:construct}
Given three lengths:
\begin{enumerate}
  \item compare the largest with the \omterm{def:g1:addition:def}{sum} of the two others; if it is
        not strictly smaller, stop: no triangle;
  \item otherwise construct as in \cref{met:g6:shapes:construct}
        (base \omterm{def:g6:lines:objects}{segment}, two compass arcs);
  \item after any construction from angles, check mentally that the
        given angles sum to less than $180^\circ$ --- with the third
        angle making exactly $180^\circ$.
\end{enumerate}
\end{method}

\section{Exercises}

\begin{exercise}[$\star$]\label{exo:g7:triangles:1}
Give the complement (to $90^\circ$) and the supplement (to
$180^\circ$) of: $30^\circ$; \ $45^\circ$; \ $72^\circ$.
\end{exercise}

\begin{exercise}[$\star$]\label{exo:g7:triangles:2}
Two lines cross; one of the four angles measures $35^\circ$. Give the
measures of the three others, with a reason for each.
\end{exercise}

\begin{exercise}[$\star$]\label{exo:g7:triangles:3}
Compute the third angle of a triangle whose first two angles measure:
(a) $40^\circ$ and $60^\circ$; \
(b) $90^\circ$ and $28^\circ$; \
(c) $75^\circ$ and $75^\circ$.
\end{exercise}

\begin{exercise}[$\star$]\label{exo:g7:triangles:4}
An \omterm{def:g6:shapes:triangles}{isosceles triangle} has its apex angle equal to $40^\circ$. Compute
its two base angles. Another \omterm{def:g6:shapes:triangles}{isosceles triangle} has a base angle of
$70^\circ$: compute its apex angle.
\end{exercise}

\begin{exercise}[$\star$]\label{exo:g7:triangles:5}
Which of these triples can be the sides of a triangle? Justify by the
triangle inequality (largest side!):
\[
(6,\ 8,\ 12); \qquad
(5,\ 5,\ 11); \qquad
(4,\ 9,\ 13); \qquad
(7,\ 7,\ 7).
\]
\end{exercise}

\begin{exercise}[$\star$]\label{exo:g7:triangles:6}
A triangle has angles $x$, $2x$ and $3x$. Find $x$ and the three
angles. What kind of triangle is it?
\end{exercise}

\begin{exercise}[$\star$]\label{exo:g7:triangles:7}
Construct a triangle $ABC$ with $BC = 6$ cm,
$\widehat{ABC} = 50^\circ$ and $\widehat{ACB} = 60^\circ$ (protractor
at $B$ and at $C$). Before constructing, predict the measure of
$\widehat{BAC}$, then check on your figure.
\end{exercise}

\begin{exercise}[$\star\star$]\label{exo:g7:triangles:8}
Two \omterm{def:g6:lines:perp}{parallel lines} are cut by a third line, making an angle of
$54^\circ$ at the first crossing (between the crossing line and one
\omterm{def:g6:lines:perp}{parallel}). Draw the situation and give the measures of all eight
angles formed.
\end{exercise}

\begin{exercise}[$\star\star$]\label{exo:g7:triangles:9}
Two sides of a triangle measure $8$ cm and $3$ cm. Between which two
values must the third side lie? Give a whole-number length that works
and one that does not.
\end{exercise}

\begin{exercise}[$\star\star$]\label{exo:g7:triangles:10}
In a triangle $ABC$, $\widehat A = 2\widehat B$ and
$\widehat C = 90^\circ$. Compute $\widehat A$ and $\widehat B$.
\end{exercise}

\begin{exercise}[$\star\star\star$]\label{exo:g7:triangles:11}
The three angles of any \omterm{def:g4:geometry:polygon}{quadrilateral} $ABCD$ can be studied by
cutting it along a diagonal into two triangles. Use this to prove that
the four angles of every \omterm{def:g4:geometry:polygon}{quadrilateral} add up to $360^\circ$, and
deduce the angle \omterm{def:g1:addition:def}{sum} of a \omterm{def:g4:geometry:polygon}{pentagon} (five sides).
\end{exercise}

\section{Problem: Walking around the block}

\begin{problem}[{Weekend problem --- the angles of any polygon:
$(n-2) \times 180^\circ$ inside, always $360^\circ$ of turning
outside, and why only three regular shapes tile a floor}]
\label{pb:g7:triangles:1}
\Cref{exo:g7:triangles:11} cut a \omterm{def:g4:geometry:polygon}{quadrilateral} into two triangles
and found $360^\circ$. This problem pushes the idea to \omterm{def:g4:geometry:polygon}{polygons}
with any number of sides, then discovers a second, completely
different proof --- by \emph{walking} around the shape and adding
up the turns --- and ends on the floor of your bathroom: among all
regular \omterm{def:g4:geometry:polygon}{polygons}, exactly three can tile a plane without gaps or
overlaps, and the bees have chosen theirs.

\medskip
\noindent\textbf{Part I --- The angles inside.}
\begin{enumerate}
  \item From one \omterm{def:g5:solids:def}{vertex} of a \omterm{def:g4:geometry:polygon}{hexagon} (six sides), draw all the
        diagonals leaving that \omterm{def:g5:solids:def}{vertex}. Into how many triangles is
        the \omterm{def:g4:geometry:polygon}{hexagon} cut, and what is the \omterm{def:g1:addition:def}{sum} of \emph{all} its
        interior angles (\cref{thm:g7:triangles:sum})?
  \item Generalize: from one \omterm{def:g5:solids:def}{vertex} of a \omterm{def:g4:geometry:polygon}{polygon} with $n$ sides,
        how many triangles does the fan of diagonals produce?
        Deduce the formula for the \omterm{def:g1:addition:def}{sum} of the interior angles,
        and check it for $n = 3$, $4$, $5$, $6$.
  \item Compute the \omterm{def:g1:addition:def}{sum} of the interior angles of a decagon
        ($10$ sides) and of a dodecagon ($12$ sides).
  \item In a \emph{regular} \omterm{def:g4:geometry:polygon}{polygon} all interior angles are
        equal. Compute the interior angle of the \omterm{def:g6:shapes:triangles}{equilateral
        triangle}, the square, the regular \omterm{def:g4:geometry:polygon}{pentagon}, \omterm{def:g4:geometry:polygon}{hexagon} and
        \omterm{def:g4:geometry:polygon}{octagon}.
  \item Explain, without any formula, why the angle \omterm{def:g1:addition:def}{sum} grows by
        exactly $180^\circ$ each time the \omterm{def:g4:geometry:polygon}{polygon} gains one
        \omterm{def:g5:solids:def}{vertex}.
\end{enumerate}

\medskip
\noindent\textbf{Part II --- The walk around the block.}
Walk along the boundary of a polygonal block, always forward. At
each corner you \emph{turn} by some angle --- the \emph{exterior
angle} of that corner; turning continues until you are back at
your starting point, facing your starting direction.
\begin{enumerate}[resume]
  \item At a corner whose interior angle is $\widehat A$, by how
        much do you turn? (Interior and exterior angle are
        supplementary.) Compute
        the three turns for a triangle with angles $40^\circ$,
        $60^\circ$, $80^\circ$.
  \item Add up the three turns of question 6. What full-circle
        fact do you observe?
  \item Explain why the turns of a complete walk around
        \emph{any} \omterm{def:g4:geometry:polygon}{polygon} --- three sides or thirty --- always
        total exactly $360^\circ$. (What has your nose done, all
        turns combined, when you arrive back?)
  \item Deduce the interior-angle formula a second time: at each
        of the $n$ corners, interior $+$ turn $= 180^\circ$; \omterm{def:g1:addition:def}{sum}
        over all corners and use question 8.
  \item In a regular \omterm{def:g4:geometry:polygon}{polygon}, all the turns are equal, so each
        exterior angle is $\frac{360^\circ}{n}$. Use this to
        answer instantly: which regular \omterm{def:g4:geometry:polygon}{polygon} has interior
        angles of $150^\circ$? And why does \emph{no} regular
        \omterm{def:g4:geometry:polygon}{polygon} have interior angles of $155^\circ$?
\end{enumerate}

\medskip
\noindent\textbf{Part III --- Tiling the floor.}
Around every point of a tiled floor, the corners of the meeting
tiles must total exactly $360^\circ$ --- no gap, no overlap
(\cref{pb:g6:lines:1}).
\begin{enumerate}[resume]
  \item For each of the \omterm{def:g6:shapes:triangles}{equilateral triangle}, the square and the
        regular \omterm{def:g4:geometry:polygon}{hexagon}: how many copies meet at a corner point
        of the tiling? Verify the $360^\circ$ each time.
  \item Show that regular \omterm{def:g4:geometry:polygon}{pentagons} cannot tile the floor: what
        do three corners total, and what would four total?
  \item Show that no regular \omterm{def:g4:geometry:polygon}{polygon} with \emph{more} than six
        sides can tile: at least three tiles must meet at each
        corner (each angle is less than $180^\circ$), and what
        happens to three corners of $135^\circ$ or more? Conclude
        the complete list of regular tilers.
  \item Bathroom floors often mix regular \omterm{def:g4:geometry:polygon}{octagons} with small
        squares. Verify that this corner works:
        $135^\circ + 135^\circ + 90^\circ$. Another classic mixes
        a square, a regular \omterm{def:g4:geometry:polygon}{hexagon} and a regular dodecagon at
        each corner: verify it too (the dodecagon's interior
        angle follows from question 10's method).
  \item The bees' choice: honeycomb cells are regular \omterm{def:g4:geometry:polygon}{hexagons}.
        Verify their corners ($3$ angles), and explain in one
        sentence --- with the help of Dido's discovery
        (\cref{pb:g6:measure:1}) --- why, of the three possible
        regular tiles, the \omterm{def:g4:geometry:polygon}{hexagon} is the cheapest in wax for
        the honey it holds.
\end{enumerate}
\end{problem}
